\section{Conclusões e perspectivas de trabalhos futuros}\label{sec:conclusao}

Este trabalho apresentou uma avaliação comparativa de técnicas de extração de características de comprimento fixo baseadas em aprendizado profundo para verificação biométrica por impressões digitais, contrastando as arquiteturas \textit{DeepPrint} e \textit{FLARE} nas edições FVC2000, FVC2002, FVC2004 e nas bases ópticas da Neurotechnology. Os resultados comprovaram a superioridade das representações densas tridimensionais frente a vetores unidimensionais: o FLARE com alinhamento por votação alcançou uma taxa média geral de EER de 0,64\%, contra 5,92\% do DeepPrint. Esse ganho decorre da preservação da topologia local das cristas e da supressão explícita de ruídos de fundo por meio de máscaras de segmentação na representação 3D, contrastando com as perdas de informação espacial inerentes à compressão em um vetor global.

A análise dos módulos revelou que a estimativa de alinhamento por votação densa superou consistentemente a predição direta por regressão no FLARE, conferindo maior estabilidade frente a distorções e impressões parciais. Quanto ao realce de imagem, a rede UNetEnh atuou de forma consistente e segura, preservando os níveis de EER das imagens brutas em ambos os modelos. Em contrapartida, a rede generativa PriorEnh induziu degradações severas em bases com forte ruído de fundo (como no FVC2000 DB1-A, no qual o EER subiu de 3,36\% para 12,45\% no DeepPrint e de 0,81\% para 8,02\% no FLARE), demonstrando que a reconstrução sintética de cristas pode introduzir artefatos e minúcias espúrias, embora em sensores ópticos de alta qualidade ambas as abordagens tenham mantido elevada fidelidade.

Como perspectivas de trabalhos futuros, destaca-se a investigação de abordagens conjuntas (\textit{end-to-end}) que treinem simultaneamente as etapas de realce, alinhamento e extração densa, mitigando a propagação de erros entre módulos independentes.
